%==============================================================================
%  CAPÍTULO IX
%==============================================================================
\chapter{Régimen especial de acciones revocatorias concursales}
\label{cap:revocatorias}

\fichaCapitulo{Ineficacia de actos lesivos a la masa de acreedores.}{Período sospechoso,
acciones objetivas y subjetivas, y exclusión del abandono del procedimiento.}

%==============================================================================
\bloqueTeorico
%==============================================================================

\subsection{Reconstrucción del patrimonio de afectación}
\pendiente{Desarrollar: protección del principio de universalidad objetiva frente a
actos fraudulentos o de mera liberalidad anteriores a la apertura.}

\subsection{Certeza jurídica del tercero de buena fe frente a la equidad de la masa}
\pendiente{Desarrollar la tensión dogmática.}

%==============================================================================
\bloqueNormativo
%==============================================================================

\subsection{Período sospechoso y plazo de interposición}

Distinción conceptual entre el \periodosospechoso{} ---época retrospectiva de celebración
del acto impugnable--- y el plazo fatal para deducir la demanda.

\subsection{La acción revocatoria objetiva}

Supuestos de los \art{287} y \art{290} \parencite{ley-20720}, en que el perjuicio a la
masa se presume
legalmente: período sospechoso de un año (pagos anticipados, daciones en pago inusuales,
garantías sobre obligaciones preexistentes) o de dos años (actos gratuitos y operaciones
con personas relacionadas).

\subsection{La acción revocatoria subjetiva}

Régimen del \art{288}: exigible sólo respecto de \empresadeudora{}s, período sospechoso
de dos años y acreditación de la \lat{participatio fraudis} del tercero. Sobre los
límites normativos de la revocación objetiva y su elusión mediante figuras afines, véase
\textcite{doc-revocacion-objetiva}; para un tratamiento monográfico,
\textcite{doc-acciones-revocatorias}.

\subsection{Plazo del artículo 291}

Un año contado desde la dictación de la resolución concursal respectiva.

%==============================================================================
\bloquePractico
%==============================================================================

\subsection{Matriz comparativa de las acciones de ineficacia}

\begin{table}[htbp]
\centering
\small
\caption{Acciones revocatorias concursales y acción pauliana civil}
\label{tab:cap09-matriz}
\begin{tabularx}{\textwidth}{@{}>{\raggedright\arraybackslash}p{0.17\textwidth} XXX@{}}
\toprule
\textbf{Elemento} & \textbf{Revocatoria objetiva (\art{287})}
& \textbf{Revocatoria subjetiva (\art{288})}
& \textbf{Pauliana civil (\artcc{2468})} \\
\midrule
Sujeto pasivo & Empresa y Persona Deudora
& Exclusivamente Empresa Deudora
& Persona Deudora, por reenvío del \art{290} \\
\addlinespace
Prueba de mala fe & No se requiere
& Exige \lat{participatio fraudis} del tercero
& Exige dolo del deudor y del tercero \\
\addlinespace
Período sospechoso & 1 año general; 2 años para actos gratuitos y relacionados
& 2 años retrospectivos
& 1 año desde la celebración del acto \\
\addlinespace
Vía procesal & Juicio sumario ante el tribunal concursal
& Juicio sumario ante el tribunal concursal
& Juicio ordinario de competencia civil común \\
\addlinespace
Abandono del procedimiento & Inaplicable
& Inaplicable
& Plenamente aplicable conforme al \cpc{} \\
\bottomrule
\end{tabularx}
\end{table}

\verificar{Contrastar cada celda con el texto vigente; en particular el reenvío del art.
290 y la exclusión del abandono, que es doctrina jurisprudencial y no regla expresa.}

\subsection{Estrategia probatoria de la acción subjetiva}
\pendiente{Desarrollar: prueba indiciaria de la participatio fraudis.}

%==============================================================================
\bloqueJurisprudencial
%==============================================================================

\subsection{Naturaleza del plazo del artículo 291: ¿prescripción o caducidad?}

El debate central consiste en determinar si el plazo de un año constituye un término de
prescripción extintiva de corto tiempo ---susceptible de interrupción civil mediante la
sola presentación de la demanda--- o un plazo de caducidad procesal fatal e improrrogable
que exige la notificación de la demanda dentro del año
\parencite{doc-plazos-revocatorias}.
\verificar{Individualizar los fallos de la \ca{Puerto Montt} y de la Sala Civil de la
\csup{} que sostienen cada posición.}

\begin{alerta}[Consecuencia práctica]
Mientras la controversia subsista, la conducta profesional prudente consiste en notificar
la demanda dentro del año, y no meramente presentarla.
\end{alerta}
